% Part: incompleteness
% Chapter: representability-in-q
% Section: beta-function

\documentclass[../../../include/open-logic-section]{subfiles}

\begin{document}

\olfileid{inc}{req}{bet}
\olsection{বিটা অপেক্ষকের সহায়ক উপপাদ্য}


যোগ, গুণ ও $\Char{=}$ পাওয়া গেলে আদিম পুনরাবৃত্তি চালানো যায়—এটি দেখাতে
আমাদের অনুক্রম সামলানোর অপেক্ষক নির্মাণ করতে হবে। (ঘাত অপেক্ষকও হাতে থাকলে
কাজটি সহজ হতো।) আদিম পুনরাবৃত্তি হাতে থাকার সময় আমরা ``$n$-তম মৌলিক
সংখ্যা''-র মতো জিনিস সংজ্ঞায়িত করে বেশ সরল একটি সংকেতায়ন বেছে নিতে
পেরেছিলাম। কিন্তু এখানে আদিম পুনরাবৃত্তি নেই—বস্তুত আমরা ন্যূনীকরণ দিয়ে
আদিম পুনরাবৃত্তি করাই দেখাতে চাই—তাই একটু বেশি কৌশলী হতে হবে।

\begin{lem}
\ollabel{lem:beta}
এমন একটি অপেক্ষক $\beta(d,i)$ আছে যাতে প্রত্যেক অনুক্রম $a_0$,
\dots,~$a_n$-এর জন্য এমন একটি সংখ্যা~$d$ আছে যে প্রত্যেক $i \le n$-এর জন্য
$\beta(d,i) = a_i$। তদুপরি, শুধু মিশ্রণ ও নিয়মিত ন্যূনীকরণ ব্যবহার করে
মৌলিক অপেক্ষকগুলি থেকে $\beta$-কে সংজ্ঞায়িত করা যায়।
\end{lem}

$d$-কে $\tuple{a_0, \dots, a_n}$ অনুক্রমটির কোড এবং $\beta(d,i)$-কে তার
$i$-তম উপাদান ফেরত দেওয়া অপেক্ষক ভাবো। (লক্ষ করো, এই ``সংকেতায়ন'' আমাদের
পরিচিত মৌলিক সংখ্যার ঘাতভিত্তিক সংকেতায়ন \emph{ব্যবহার করে না}!) সহায়ক
উপপাদ্যটির দাবি খুবই সামান্য; অনুক্রম জোড়া লাগানো বা শেষে উপাদান যোগ করা
যায়, এমনকি মিশ্রণ ও নিয়মিত ন্যূনীকরণে সংজ্ঞেয় অপেক্ষক ব্যবহার করে $a_0$,
\dots,~$a_n$ থেকে $d$ \emph{গণনা} করা যায়—এসবের কোনোটিই এটি বলে না। শুধু
বলে, এমন একটি ``অসংকেতায়ন'' অপেক্ষক আছে যার দ্বারা প্রত্যেক অনুক্রম
``সংকেতায়িত'' হয়।

$\beta$ প্রতীকটি গ্যোডেলের। আবার বলি, সহায়ক উপপাদ্য প্রমাণের কঠিন অংশ হলো
আপাতদৃষ্টিতে সীমিত উপায়—শুধু মিশ্রণ ও ন্যূনীকরণ—দিয়ে উপযুক্ত $\beta$
সংজ্ঞায়িত করা; তবে যোগ, গুণ ও~$\Char{=}$ ব্যবহার করা যাবে। এই সহায়ক উপপাদ্য
প্রমাণের নানা উপায় আছে; সবচেয়ে পরিচ্ছন্ন উপায়গুলির একটি এখনও গ্যোডেলের
মূল পদ্ধতি। এতে সুন্ৎসির উপপাদ্য নামে একটি সংখ্যাতাত্ত্বিক ফল ব্যবহৃত হয়
(প্রথাগতভাবে ``চীনা অবশিষ্ট উপপাদ্য'')।

\begin{defn}
দুটি স্বাভাবিক সংখ্যা $a$ ও $b$-র গরিষ্ঠ সাধারণ গুণনীয়ক~$1$ হলে এবং কেবল
তখনই তাদের \emph{পরস্পর সহমৌলিক} বলা হয়; অন্যভাবে বললে, ১ ছাড়া তাদের
আর কোনো সাধারণ গুণনীয়ক নেই।
\end{defn}

\begin{defn}
স্বাভাবিক সংখ্যা $a$ ও $b$-কে \emph{মডুলো~$c$ অনুসারে সর্বসম} বলা হয়,
$a \equiv b \mod c$, যদি এবং কেবল যদি $c \mid (a-b)$; অর্থাৎ $c$ দিয়ে ভাগ
করলে $a$ ও $b$-র একই ভাগশেষ থাকে।
\end{defn}

সুন্ৎসির উপপাদ্যটি হলো:
\begin{thm}
ধরি $x_0$, \dots,~$x_n$ (জোড়ায় জোড়ায়) পরস্পর সহমৌলিক। $y_0$,
\dots,~$y_n$ যেকোনো সংখ্যা হোক। তাহলে এমন একটি সংখ্যা $z$ আছে যাতে
\begin{align*}
z & \equiv y_0 \mod x_0 \\
z & \equiv y_1 \mod x_1 \\
& \vdots  \\
z & \equiv y_n \mod x_n.
\end{align*}
\end{thm}

সুন্ৎসির উপপাদ্য আমরা এভাবে ব্যবহার করব: $x_0$, \dots,~$x_n$ যথাক্রমে
$y_0$, \dots,~$y_n$-এর চেয়ে বড় হলে $z$-কে
$\tuple{y_0, \dots,y_n}$ অনুক্রমটির কোড নেওয়া যায়। $y_i$ ফিরে পেতে শুধু
$z$-কে $x_i$ দিয়ে ভাগ করে ভাগশেষ নিতে হবে। এই সংকেতায়ন ব্যবহার করতে
$x_0$, \dots,~$x_n$-এর উপযুক্ত মান বের করতে হবে।

কয়েকটি পর্যবেক্ষণ এ কাজে সাহায্য করবে। $y_0$, \dots,~$y_n$ দেওয়া থাকলে
ধরি
\begin{align*}
j &= \max(n, y_0 + 1, \dots, y_n + 1), \\
m &= \lcm(1,\dots,j),
\end{align*}
এবং ধরি
\begin{align*}
x_0 & = 1 + m \\
x_1 & = 1 + 2 \cdot m \\
x_2 & = 1 + 3 \cdot m \\
& \vdots  \\
x_n & = 1 + (n+1) \cdot m
\end{align*}
তাহলে দুটি কথা সত্য:
\begin{enumerate}
\item\ollabel{rel-prime} $x_0,\dots,x_n$ পরস্পর সহমৌলিক।
\item\ollabel{less} প্রত্যেক $i$-এর জন্য $y_i < x_i$।
\end{enumerate}
\olref{rel-prime} সত্য দেখতে লক্ষ করো, $p$ মৌলিক সংখ্যা এবং $p \mid x_i$
ও $p \mid x_k$ হলে $p \mid 1 + (i+1) m$ ও
$p \mid 1 + (k+1) m$। তখন তাদের অন্তরকেও $p$ ভাগ করে,
\[
(1 + (i+1)m) - (1+ (k+1)m) = (i-k) m.
\]
যেহেতু $p$ সংখ্যা $1 + (i+1)m$-কে ভাগ করে, তাই একই সঙ্গে $m$-কে ভাগ করতে
পারে না (তা করলে প্রথম ভাগে ভাগশেষ~$1$ থাকত)। অতএব $p$ সংখ্যা
$(i-k)m$-কে $p$ ভাগ করে বলে $i-k$-কেও ভাগ করে। কিন্তু
$\left|i-k\right|$ সর্বাধিক~$n$, আর আমরা $j \geq n$ বেছেছি; তাই এর থেকে
$p \mid m$ হয়—আবার বিরোধ। সুতরাং $x_i$ ও $x_k$ দুটিকেই ভাগ করে এমন মৌলিক
সংখ্যা নেই। অনুচ্ছেদ~\olref{less} সহজ:
$y_i < j \leq m < x_i$।

এবার $\beta$ অপেক্ষকের সহায়ক উপপাদ্য প্রমাণ করি। মনে রাখো, আমরা $0$,
উত্তরসূরি, যোগ, গুণ, $\Char{=}$, অভিক্ষেপ এবং এগুলি থেকে নিয়মিত অপেক্ষকের
উপর মিশ্রণ ও ন্যূনীকরণ ব্যবহার করে সংজ্ঞায়িত যেকোনো অপেক্ষক নিতে পারি।
কোনো সম্বন্ধের চরিত্রসূচক অপেক্ষক এভাবে সংজ্ঞেয় হলে সেই সম্বন্ধটিও ব্যবহার
করা যায়। আগের মতোই দেখানো যায় যে এই সম্বন্ধগুলি বুলীয় মিশ্রণ ও সীমাবদ্ধ
পরিমাণায়নের অধীনে বদ্ধ; যেমন:
\begin{align*}
\fn{not}(x) & \defis \Char{=}(x,0)\\
\bmin{x \leq z}{R(x,y)} & \defis \umin{x}{(R(x,y) \lor x = z)}\\
\bexists{x \leq z}{R(x,y)} & \defiff R(\bmin{x \leq z}{R(x,y)}, y)
\end{align*}
এরপর দেখানো যায় যে নিচের সবকটিও আদিম পুনরাবৃত্তি ছাড়াই সংজ্ঞেয়:
\begin{enumerate}
\item জোড়-গঠন অপেক্ষক, $J(x,y) = \frac{1}{2}[(x+y)(x+y+1)] + x$;
% maybe explain more what is going on here, a bit confusing.
\item অভিক্ষেপ অপেক্ষকদুটি
\begin{align*}
K(z) & = \bmin{x \leq z}{\bexists{y \leq z}{z = J(x,y)}},\\
L(z) & = \bmin{y \leq z}{\bexists{x \leq z}{z = J(x,y)}};
\end{align*}
\item ক্ষুদ্রতর-সম্বন্ধ $x < y$;
\item বিভাজ্যতা-সম্বন্ধ $x \mid y$;
% \item $x \tsub y$
% \item $\fn{Prime}(x)$
% \item Assuming $p$ is prime, the relation ``$x$ is a power of $p$'':
% \[
% \bforall{y \leq x}{(y \mid x \lif y = 1 \lor y = x)}.
% \]
\item অপেক্ষক $\fn{rem}(x,y)$, যা $y$-কে $x$ দিয়ে ভাগ করলে ভাগশেষটি
  ফেরত দেয়।
\end{enumerate}
এবার সংজ্ঞায়িত করি
\begin{align*}
\beta^*(d_0,d_1,i) & = \fn{rem}(1+(i+1) d_1,d_0) \text{ এবং}\\
\beta(d,i) & = \beta^*(K(d),L(d),i).
\end{align*}
এটিই আমাদের চাওয়া অপেক্ষক। উপরের মতো $a_0,\dots,a_n$ দেওয়া থাকলে ধরি
\[
j = \max(n,a_0+1,\dots,a_n+1),
\]
এবং $d_1 = \lcm(1,\dots,j)$ নিই। উপরের \olref{rel-prime} থেকে জানি,
$1+d_1$, $1+2 d_1$, \dots, $1+(n+1) d_1$ পরস্পর সহমৌলিক; আর
\olref{less} থেকে জানি, এগুলি যথাক্রমে $a_0,\dots,a_n$-এর চেয়ে বড়।
সুন্ৎসির উপপাদ্য অনুসারে এমন একটি মান~$d_0$ আছে যাতে প্রত্যেক~$i$-এর জন্য
\[
d_0 \equiv a_i \mod (1+(i+1)d_1)
\]
এবং তাই ($d_1$ যেহেতু $a_i$-এর চেয়ে বড়),
\[
a_i = \fn{rem}(1+(i+1)d_1,d_0).
\]
$d = J(d_0,d_1)$ ধরি। তাহলে প্রত্যেক $i \le n$-এর জন্য
\begin{align*}
\beta(d,i) & =  \beta^*(d_0,d_1,i) \\
& =  \fn{rem}(1+(i+1) d_1,d_0) \\
& =  a_i
\end{align*}
পাই, আর এটিই আমাদের দরকার। এতে $\beta$-অপেক্ষকের সহায়ক উপপাদ্যের প্রমাণ
সম্পূর্ণ হলো।

\begin{prob}
  দেখাও, $x < y$, $x \mid y$ সম্বন্ধদুটি এবং অপেক্ষক~$\fn{rem}(x,y)$-কে
  আদিম পুনরাবৃত্তি ছাড়াই সংজ্ঞায়িত করা যায়। $0$, উত্তরসূরি, যোগ, গুণ,
  $\Char{=}$, অভিক্ষেপ এবং সীমাবদ্ধ ন্যূনীকরণ ও পরিমাণায়ন ব্যবহার করতে পারো।
\end{prob}

\end{document}
